% =========================================================================
%  Ejercicio 1f - Redundancia de datos 
%  Responsable: [León Bautista Yahir]
% =========================================================================

\section*{Ejercicio 1f. Redundancia de datos}
Investiga qué es la redundancia de datos. ¿Cuál sería la diferencia entre redundancia de datos
controlada y no controlada? Proporciona un ejemplo claro de cada tipo. ¿Por qué la redundancia
no controlada suele ser un problema para la integridad y consistencia de los datos?

La redundancia de datos ocurre cuando los mismos datos se almacenan en múltiples copias, ubicaciones o sistemas diferentes.
\begin{itemize}
    \item Redundancia de datos controlada: Se refiere a la duplicación de datos que se gestiona de manera intencional y planificada, 
    generalmente para mejorar el rendimiento o la disponibilidad del sistema. Un ejemplo de redundancia controlada es el uso de bases de datos replicadas
    en diferentes servidores para garantizar la disponibilidad continua de los datos en caso de fallos del sistema.
    \item Redundancia de datos no controlada: Se refiere a la duplicación de datos que ocurre de manera accidental o sin una gestión adecuada, 
    lo que puede llevar a inconsistencias y errores en los datos. Un ejemplo de redundancia no controlada es cuando un usuario introduce la misma 
    información en diferentes formularios o sistemas sin coordinación, lo que puede resultar en registros duplicados y conflictos de datos cuando se realicen modificaciones.
\end{itemize}
La redundancia no controlada suele ser un problema para la integridad y consistencia de los datos porque puede generar conflictos y errores en la información almacenada.
Cuando un dato existe en varios lugares sin supervisión del sistema, inevitablemente surgen problemas conocidos como anomalías de modificación, que se generan al actualizar, insertar o eliminar datos en diferentes ubicaciones sin coordinación.


% Ref: Elmasri, R., & Navathe, S. B. (2016). Fundamentals of database systems (7th ed.). Pearson. 

% Responder aqui. Contenido minimo:
% - Definir que es la redundancia de datos (investigacion con fuentes).
% - Diferencia entre redundancia CONTROLADA y NO controlada.
% - Ejemplo claro de cada tipo.
% - Por que la redundancia no controlada suele ser un problema para la
%   integridad y consistencia de los datos.